\scanpage{153}
% 153.01 Top centre
\mspage{476}

% 153.02 Top paragraph, note 108 with following 1
(\textsuperscript{108}1) (6 octobre) Je veux dire que les forces de nature répressive qui ont joué dans ma vie, semblent prendre surtout, sinon exclusivement, l'une de ces deux formes spécifiques: enterrement du passé, et mise en avant de mes traits ``virils au détriment de mes traits ``féminins''. Je n'entends nullement dire que ces deux ``forces'', de nature répressive l'une et l'autre (c'est à dire, visant à un ``refoulement'', à un escamotage d'une certaine réalité), soient les seules qui aient ``dominé ma vie''! Ce serait oublier tout l'aspect non égotique de mon être, la pulsion de connaissance s'exprimant aussi bien au niveau de corps que de \unsure{l'espri|l'esprit} (Voir à ce sujet notamment ``Mes passions'', section n\textsuperscript{o} 35.)

% 153.03 Middle paragraph
Même parmi les forces stucturant le moi, émanation du ``patron'' donc, il en est au moins une, de nature non répressive par elle-même, bien antérieure aux forces de répression et dont le rôle dans ma vie a été plus essentiel encore: c'est l'identification à mon père, qui a été comme ``le coeur paisible et \unsure{puissan|puissant} du sentiment de ma propre force. Cette identification n'allait nullement dans le sens de l'exaltation de certaines valeurs ou qualités (viriles disons) au détriment d'autres (``féminines''). Indépendamment des valeurs professées par mon père, sa personne (jusqu'en 1933, où un basculement a eu lieu en lui (*)), a été empreinte d'un fort équilibre yin-yang, où l'intuition et la spontaneité n'avaient pas une moindre part que l'intellect et la volonté.

% 153.04 Lower paragraph
Enfin, comme autre ``force'' importante de nature égotique, intimement liée, elle, aux mécanismes répressifs (ou pour mieux dire, de nature ``répressive'' elle-même), il convient encore de compter la sempiternelle v\,a\,n\,i\,t\,é , dont le rôle a été aussi lourd dans ma vie que dans celle de quiconque. Mais cette ``force''-là est de nature si universelle, tout comme le rôle dominant qu'elle joue dans la vie de chacun (sous une forme plus ou moins grossière ou subtile), qu'il n'y a guère lieu de l'inclure expressément, dans un relevé des formes spécifiques que prennent chez une personnne les forces et mécanismes qui structurent le moi, et donnent à celui-ci sa physionomie particulière et son assise.

% 153.05 Footnote below rule
\noindent\rule{5em}{0.3pt}\par
(*) Chose remarquable\handwritten{,} ce ``basculement'' chez mon père (agé alors de 43 ans) s'est fait vers un état s\,u\,p\,e\,r\,-\,y\,i\,n , vers une sorte de passivité de pacha, en étroite connivence avec ma mère, jouant un rôle super-yang. Celle-ci l'a pris en charge \hspace{1em} lieu et place de ses enfants.(Ils sont largués aux ``profits et pertes'', jusqu'en 1939 tout au moins, l'année où sous la pression des évènements et à son corps défendant, elle finit par me reprendre auprès d'elle...) Cette relation de dépendance de mon père et de renversement des rôles yin-yang entre mes parents, a duré jusqu'à la disparition de mon père en 1942.
% archive sheet number: 152
